\documentclass[../../main.tex]{subfiles}

\begin{document}

\chapter{Kriptografi Klasik}
\label{chap:kriptografi-klasik}

\begin{learningoutcome}
Setelah mempelajari bab ini, mahasiswa diharapkan mampu:
\begin{itemize}
    \item Menjelaskan prinsip dasar cipher substitusi dan transposisi serta memberikan contoh implementasinya (Sub-CPMK 1.1).
    \item Menganalisis kelemahan cipher abjad-tunggal dan menerapkan teknik analisis frekuensi untuk memecahkan cipherteks (Sub-CPMK 1.1).
    \item Mengimplementasikan Vigen\`ere Cipher dan memahami konsep cipher abjad-majemuk (\textit{polyalphabetic}) (Sub-CPMK 1.1).
    \item Menjelaskan mekanisme kerja Playfair Cipher sebagai bentuk substitusi digram (Sub-CPMK 1.1).
    \item Menerapkan konsep Affine Cipher dan Hill Cipher berbasis aljabar linier (Sub-CPMK 1.1).
    \item Memahami sejarah dan mekanisme kerja mesin Enigma sebagai tonggak kriptografi elektromekanik (Sub-CPMK 1.1).
    \item Menjelaskan prinsip \textit{One-Time Pad} beserta syarat kerahasiaan sempurna dan kelemahan praktisnya (Sub-CPMK 1.1).
\end{itemize}
\end{learningoutcome}

\subfile{section-01}
\subfile{section-02}
\subfile{section-03}
\subfile{section-04}
\subfile{section-05}
\subfile{section-06}
\section{Contoh Terhitung}

\begin{example}
\textbf{Vigen\`ere Cipher.} Enkripsi plainteks \code{SERANGPAGI} dengan kunci \code{KUNCI}. Kunci diulang periodik sepanjang pesan sehingga menjadi \code{KUNCIKUNCI} ($|P|=|K|=10$). Konversi huruf ke angka ($A=0,\dots,Z=25$):
\[
\begin{array}{c|cccccccccc}
P & 18 & 4 & 17 & 0 & 13 & 6 & 15 & 0 & 6 & 8\\
K & 10 & 20 & 13 & 2 & 8 & 10 & 20 & 13 & 2 & 8
\end{array}
\]
Hitung $C_i=(P_i+K_i)\bmod 26$ tiap posisi:
\[
\begin{aligned}
(18+10)_{26}&=2=\text{C}, & (13+8)_{26}&=21=\text{V},\\
(4+20)_{26}&=24=\text{Y}, & (6+10)_{26}&=16=\text{Q},\\
(17+13)_{26}&=4=\text{E}, & (15+20)_{26}&=9=\text{J},\\
(0+2)_{26}&=2=\text{C}, & (0+13)_{26}&=13=\text{N},\\
(6+2)_{26}&=8=\text{I}, & (8+8)_{26}&=16=\text{Q}.
\end{aligned}
\]
Jadi cipherteks adalah \code{CYECVQJNIQ}. Dekripsi mengembalikannya melalui $P_i=(C_i-K_i)\bmod 26$.
\end{example}

\begin{example}
\textbf{Hill Cipher $2\times2$.} Matriks kunci $\mathbf{K}=\begin{pmatrix}3&3\\2&5\end{pmatrix}$ dengan $\det(\mathbf{K})=(3\cdot5)-(3\cdot2)=9$, dan $\gcd(9,26)=1$ sehingga $\mathbf{K}$ invertibel modulo 26. Enkripsi plainteks \code{HILL} sebagai dua vektor kolom: $\text{HI}=(7,8)$ dan $\text{LL}=(11,11)$.
\[
\mathbf{C}=\mathbf{KP}\pmod{26}
\]
\[
\begin{pmatrix}3&3\\2&5\end{pmatrix}\begin{pmatrix}7\\8\end{pmatrix}=
\begin{pmatrix}21+24\\14+40\end{pmatrix}=\begin{pmatrix}45\\54\end{pmatrix}\bmod 26=\begin{pmatrix}19\\2\end{pmatrix}=\text{TC}
\]
\[
\begin{pmatrix}3&3\\2&5\end{pmatrix}\begin{pmatrix}11\\11\end{pmatrix}=
\begin{pmatrix}33+33\\22+55\end{pmatrix}=\begin{pmatrix}66\\77\end{pmatrix}\bmod 26=\begin{pmatrix}14\\25\end{pmatrix}=\text{OZ}
\]
Cipherteks \code{TCOZ}. Dekripsi memakai $\mathbf{K}^{-1}=\det(\mathbf{K})^{-1}\begin{pmatrix}5&-3\\-2&3\end{pmatrix}$ dengan $9^{-1}\equiv3 \pmod{26}$, yaitu $\mathbf{K}^{-1}=\begin{pmatrix}15&17\\20&9\end{pmatrix}$. Uji balik untuk \code{TC}: $\mathbf{K}^{-1}\binom{19}{2}=\binom{15\cdot19+17\cdot2}{20\cdot19+9\cdot2}=\binom{319}{398}\bmod26=\binom{7}{8}=\text{HI}$.
\end{example}

\begin{example}
\textbf{One-Time Pad (XOR biner).} Plainteks $P=10110011$ dienkripsi dengan kunci acak $K=01101101$; panjang kunci sama dengan panjang pesan ($|K|=|P|=8$ bit). Enkripsi dan dekripsi identik, yaitu XOR bit-per-bit:
\[
C=P\oplus K=10110011\oplus01101101=11011110.
\]
Dekripsi: $C\oplus K=11011110\oplus01101101=10110011=P$. Jika kunci yang sama dipakai lagi untuk $P_2=11001010$, diperoleh $C_2=10100111$; maka $C\oplus C_2=01111001=P\oplus P_2$, sehingga XOR dua plainteks terbuka tanpa mengetahui kuncinya.
\end{example}


\begin{obeactivity}
\textbf{Aktivitas 4.1: Praktik Analisis Frekuensi}
1. Gunakan alat \textit{online} seperti \url{https://www.cryptool.org/en/cto/n-gram-analysis}.
2. Masukkan sebuah cipherteks yang dihasilkan dari monoalphabetic cipher.
3. Analisis distribusi frekuensi hurufnya dan coba petakan huruf yang paling sering muncul ke huruf 'E' atau 'A'.
4. Lakukan iterasi \textit{trial-and-error} hingga pesan dapat terbaca sepenuhnya.
\end{obeactivity}


\begin{obereflection}
\textbf{Latihan 4.1}
1. Gunakan kunci \texttt{Kripto} (dengan $m=7, b=10, n=26$) untuk mengenkripsi kata \texttt{halo}.
2. Diberikan cipherteks Vigen\`ere \texttt{LSimMNVWGRR} dengan kunci \texttt{soreang}, lakukan dekripsi untuk menemukan plainteksnya.
3. Buatlah matriks kunci $2 \times 2$ untuk Hill Cipher, hitung determinannya, dan tentukan apakah matriks tersebut memiliki invers modulo 26.
4. Dekripsikan cipherteks Playfair \texttt{ZB RS FY KU PG LG RK VS NL QV} dengan matriks kunci yang dibentuk dari kalimat kunci \texttt{JALAN GANESHA SEPULUH}. Tunjukkan langkah pembuangan huruf pengisi \texttt{X} pada akhir dekripsi.
5. Diketahui dua pasang plainteks--cipherteks pada affine cipher: huruf \texttt{K} menghasilkan \texttt{D}, dan huruf \texttt{T} menghasilkan \texttt{A}. Tentukan nilai $m$ dan $b$, lalu periksa apakah $\gcd(m,26)=1$.
6. Pada sebuah cipherteks Vigen\`ere ditemukan kriptogram berulang dengan jarak 15, 20, dan 30 huruf. Tentukan panjang kunci yang paling mungkin beserta alasan langkah demi langkah.
7. Enkripsikan plainteks \texttt{KRIPTOGRAFI} dengan Rail Fence cipher empat baris, lalu jelaskan mengapa cipherteks itu masih dapat dipecahkan meskipun jumlah barisnya ditambah.
\end{obereflection}

\section{Rangkuman}
\begin{itemize}
    \item Cipher klasik dibedakan atas substitusi (Caesar, monoalphabetic, Vigen\`ere, Affine) dan transposisi (kolom); keduanya bekerja pada huruf, bukan pada blok bit.
    \item Cipher abjad-tunggal mudah dipecahkan karena distribusi frekuensi huruf plainteks tetap tampak pada cipherteks (analisis frekuensi, bigram, dan trigram).
    \item Vigen\`ere dan cipher abjad-majemuk meratakan frekuensi, tetapi dapat dipecahkan dengan metode Kasiski bila pesan cukup panjang dan kunci diulang periodik.
    \item Playfair dan Hill adalah polygram cipher; Hill berbasis aljabar linier $\mathbf{C}=\mathbf{KP}\bmod 26$ dan hanya dapat didekripsi jika $\det(\mathbf{K})$ relatif prima dengan 26.
    \item Affine Cipher memakai fungsi linier $C\equiv mP+b\pmod n$ dan mensyaratkan $\gcd(m,n)=1$ agar invers modulo ada.
    \item Enigma adalah mesin elektromekanik abjad-majemuk berperiode sangat panjang; kelemahannya dimanfaatkan kriptanalisis Polandia--Inggris (mesin \textit{Bombe} Turing).
    \item One-Time Pad (Vernam 1917; dibuktikan Shannon 1949) adalah satu-satunya cipher berkerahasiaan sempurna, dengan syarat kunci acak sejati, panjang kunci sama dengan panjang pesan, dan kunci hanya dipakai sekali.
    \item Kelemahan praktis OTP: pembangkitan dan distribusi kunci sepanjang pesan sulit, serta pemakaian ulang kunci (\textit{two-time pad}) langsung membocorkan XOR dua plainteks.
\end{itemize}


\begin{obeassessment}
\textbf{Kuis Bab 4}
1. Mengapa Vigen\`ere Cipher lebih aman daripada Caesar Cipher terhadap serangan analisis frekuensi?
2. Jelaskan mengapa dalam Playfair Cipher huruf 'J' biasanya dihilangkan atau digabung dengan 'I'.
3. Apa syarat utama agar sebuah matriks dapat digunakan sebagai kunci dalam Hill Cipher?
4. Jelaskan mengapa parameter $m$ pada affine cipher harus relatif prima dengan ukuran alfabet, dan berikan satu contoh nilai $m$ yang tidak sah beserta akibatnya bagi proses dekripsi.
5. Diberikan cipherteks Vigen\`ere dengan kriptogram berulang berjarak 15, 20, dan 30 huruf. Tentukan panjang kunci yang paling mungkin dan uraikan langkah-langkahnya.
6. Sebutkan tiga syarat kerahasiaan sempurna One-Time Pad, lalu jelaskan secara singkat apa yang terjadi bila masing-masing syarat dilanggar.
\end{obeassessment}
\begin{competencychecklist}
\begin{tabularx}{\linewidth}{|X|c|}
\hline
Indikator Kompetensi & Tercapai \\
\hline
Saya dapat membedakan cipher substitusi dan transposisi & [ ] \\
\hline
Saya dapat melakukan analisis frekuensi untuk memecahkan monoalphabetic cipher & [ ] \\
\hline
Saya dapat melakukan enkripsi dan dekripsi Vigen\`ere Cipher & [ ] \\
\hline
Saya dapat menerapkan aturan enkripsi Playfair Cipher & [ ] \\
\hline
Saya dapat menghitung operasi matriks pada Hill Cipher & [ ] \\
\hline
Saya dapat menjelaskan prinsip kerja rotor pada mesin Enigma & [ ] \\
\hline
Saya dapat melakukan enkripsi dan dekripsi Affine Cipher serta menentukan ruang kuncinya & [ ] \\
\hline
Saya dapat melakukan kriptanalisis Hill dan Affine Cipher dengan serangan \textit{known-plaintext} & [ ] \\
\hline
Saya dapat menentukan panjang kunci Vigen\`ere dengan metode Kasiski & [ ] \\
\hline
Saya dapat menjelaskan kerahasiaan sempurna One-Time Pad beserta tiga syaratnya & [ ] \\
\hline
Saya dapat menjelaskan mengapa One-Time Pad tidak sesuai untuk komunikasi modern & [ ] \\
\hline
\end{tabularx}
\end{competencychecklist}


\end{document}
